\SourcePage{34}{37}
\section{The Limiting Transition}

Classical mechanics is contained in quantum mechanics as a limiting case. We
must determine how this limiting transition takes place.

In quantum mechanics a wave function describes the electron, governing the
various possible values of its coordinate; thus far we know
only that this function solves a certain linear partial differential equation.
Classical mechanics, by contrast, treats the electron as a material particle
moving along a trajectory fixed completely by the equations of motion. A
relation analogous in a certain sense to that between quantum and classical
mechanics occurs in electrodynamics between wave optics and geometrical
optics. Wave optics describes electromagnetic waves by electric- and
magnetic-field vectors satisfying a particular system of linear differential
equations (Maxwell's equations). Geometrical optics, however,%
\SourcePage{35}{38}
describes the propagation of light along definite paths---rays. This analogy
leads one to conclude that the quantum-to-classical limiting transition is
analogous to that from wave optics to geometrical optics.

We now recall how this latter transition is carried out mathematically (see II,
\secref{53}). Let $u$ be any component of the field in an electromagnetic wave. It
may be written $u=ae^{i\varphi}$, with real amplitude $a$ and phase $\varphi$
(the phase is called the eikonal in geometrical optics). The geometrical-optics
limit corresponds to short wavelengths; mathematically, this means that
$\varphi$ changes greatly over short distances. In particular, its absolute
magnitude may be regarded as large.

Correspondingly, assume that the classical limiting case in quantum mechanics
is represented by wave functions of the form $\Psi=ae^{i\varphi}$, where $a$
varies slowly and $\varphi$ takes large values. In mechanics, as is known, a
particle trajectory can be found from a variational principle according to
which the so-called action $S$ of the mechanical system must be a minimum (the
principle of least action). In geometrical optics the ray path is instead
fixed by Fermat's principle: the ``optical path length'' of the ray, that is,
the difference between its phases at the end and at the beginning of the path,
must be a minimum.

This analogy allows us to conclude that, in the classical limiting case, the
phase $\varphi$ of the wave function must be proportional to the mechanical
action $S$ of the physical system under consideration: $S=\operatorname{const}
\cdot\varphi$. The proportionality factor is called \emph{Planck's constant}
and is denoted by $\hbar$\footnote{It was introduced into physics by
\emph{Planck} (\emph{M.~Planck}, 1900). The constant $\hbar$ used throughout
this book is, strictly speaking, Planck's constant $h$ divided by $2\pi$
(Dirac's notation).}. Since $\varphi$ is dimensionless, it has the dimensions
of action, and its value is
\[
  \hbar=1{,}055\cdot10^{-27}\,\text{erg}\cdot\text{s}.
\]

Thus the wave function of an ``almost classical'' (or, as it is called,
\emph{quasiclassical}) physical system has the form
\begin{equation}
  \Psi=ae^{iS/\hbar}.
  \tag{6.1}
\end{equation}
\SourcePage{36}{39}
Planck's constant has a fundamental part in every quantum phenomenon. Its
relative magnitude (in comparison with other quantities having the same
dimensions) determines the ``degree of quantumness'' of a given physical
system. Passage from quantum to classical mechanics corresponds to a large
phase and can formally be described by taking the limit $\hbar\to0$ (just as
the passage from wave to geometrical optics corresponds to the zero-wavelength
limit, $\lambda\to0$).

We have established the limiting form of the wave function, but must still ask
how it is related to classical motion along a trajectory. In general, motion
described by a wave function does not at all turn into motion along one fixed
trajectory. Its connection with classical motion is this: if the wave function,
and hence the probability distribution of the coordinates, is specified at an
initial instant, the distribution thereafter ``moves'' as required by the
laws of classical mechanics (see the end of \secref{17} for details).

To obtain motion along a definite trajectory one must begin with a wave
function of a special kind, appreciably different from zero only within a very
small region of space (a so-called \emph{wave packet}); the dimensions of this
region may be made to tend to zero together with $\hbar$. One can then state
that, in the quasiclassical case, the wave packet moves through space along the
particle's classical trajectory.

Finally, in the limit, quantum-mechanical operators must reduce simply to
multiplication by their corresponding physical quantities.
